\documentclass[11pt]{article}

\usepackage[margin=1in]{geometry}
\usepackage{amsmath,amssymb,amsthm,mathtools}
\usepackage{booktabs}
\usepackage{microtype}
\usepackage[hidelinks]{hyperref}
\usepackage{url}

\newtheorem{theorem}{Theorem}[section]
\newtheorem{corollary}[theorem]{Corollary}
\newtheorem{proposition}[theorem]{Proposition}
\newtheorem{lemma}[theorem]{Lemma}
\theoremstyle{definition}
\newtheorem{definition}[theorem]{Definition}
\theoremstyle{remark}
\newtheorem{remark}[theorem]{Remark}

\newcommand{\Vm}{V_m}
\newcommand{\Km}{\mathcal K_m}
\newcommand{\Rm}{\mathcal R_m}
\newcommand{\R}{\mathbb R}
\newcommand{\Q}{\mathbb Q}
\newcommand{\ip}[2]{\langle #1,#2\rangle}
\newcommand{\norm}[1]{\lVert #1\rVert}
\newcommand{\gap}{\gamma_m}

\title{Weighted-Span Obstructions for Sparse Elastic Rectangle Frames on Velocity Grids}
\author{Runyuan Wang}
\date{Draft revised 21 August 2026}

\begin{document}
\maketitle

\begin{abstract}
Let $V_m=\{-m,\ldots,m\}^3\subset\mathbb Z^3$.  An elastic
rectangle gives a signed four-point row that annihilates the five collision
modes $1,x,y,z,$ and $|v|^2$.  For a family of such rows, let
$C_m$ be the signed matrix, let $D_m$ be the diagonal matrix of vertex
incidence degrees, and use the canonical normalization
$A_m=\tfrac12 C_mD_m^{-1/2}$.  We give a self-contained weighted-span
obstruction.  If $R$ has orthogonal side vectors $p_R,q_R$, then the least
singular value of $A_m$ on the quotient by the five weighted collision modes
obeys
\[
 \gamma_m^2\leq
 G_m^{-1}\sum_{R\in\mathcal R_m}(p_{R,x}q_{R,x}-p_{R,y}q_{R,y})^2
 \leq G_m^{-1}\sum_{R\in\mathcal R_m}|p_R|^2|q_R|^2,
\]
where $G_m=\sum_{v\in V_m}(x^2-y^2)^2$.  Thus uniformly bounded degree
and sublinear side lengths $L_m=o(m)$ cannot produce a uniform normalized
gap.  Conversely, any bounded-degree family with $\gamma_m\geq c>0$ must put
a positive fraction of rectangles at quantitatively macroscopic scale.  This
condition is necessary, not sufficient.

We separately report four finite full-pool residual-greedy instances for
$m=2,3,4,5$.  Exact odd-minor certificates show that each of their rational
kernels is precisely the five-dimensional collision space, while numerical
normalized gaps range from approximately $0.138$ to $0.104$.  The exact
kernel statements and the floating-point gaps are finite evidence only.
Only the quantitatively macroscopic, diameter-scale regime remains open for
the unrestricted nonlocal bounded-degree proposition.
\end{abstract}

\section{Introduction and scope}\label{sec:intro}

This paper asks whether collision-type incidence structures can support
sparse, well-conditioned constraint systems.  The result proved here is a
statement about finite matrices built from orthogonal velocity rectangles.

The paper has two deliberately separated components.  First, we prove an
explicit upper bound for the normalized quotient gap when both vertex degree
and rectangle side length are bounded.  This is an infinite-family theorem,
but only under those two hypotheses.  Second, we describe four finite
matrices selected from complete rectangle pools.  Their rational kernels are
certified exactly; their gaps are computed numerically and independently
recomputed from the raw row lists.  Four finite cases do not establish a
uniform asymptotic lower bound.

For clarity, the evidential status of assertions is as follows.
\begin{center}
\begin{tabular}{@{}clp{0.58\textwidth}@{}}
\toprule
label & status & meaning in this paper \\
\midrule
R & rigorous & derived symbolically for every object satisfying the stated hypotheses;\\
E & exact finite & certified over $\mathbb Z$ or $\mathbb Q$ for a named finite matrix;\\
N & numerical finite & floating-point computation with residual and reproducibility checks;\\
O & open & neither proved nor disproved by the available argument or data.\\
\bottomrule
\end{tabular}
\end{center}
The weighted-span obstruction is R.  The rational ranks and degree counts in
Section~\ref{sec:finite} are E.  The gaps in that section are N.  The
unrestricted nonlocal proposition in Section~\ref{sec:open} is O.

\section{Elastic rectangles and the normalized quotient}\label{sec:operator}

\subsection{Rectangle rows}

Fix an integer $m\geq1$ and write
\[
 V_m=\{-m,-m+1,\ldots,m\}^3,\qquad N_m=|V_m|=(2m+1)^3.
\]
An elastic rectangle in $V_m$ is a quadruple that can be represented as
\begin{equation}\label{eq:rectangle}
 a,\qquad a+p,\qquad a+q,\qquad a+p+q,
 \qquad p,q\in\mathbb Z^3\setminus\{0\},\quad p\cdot q=0,
\end{equation}
with all four vertices in $V_m$.  Its momentum identity is automatic.  The
energy residual equals
\[
 |a|^2+|a+p+q|^2-|a+p|^2-|a+q|^2=2p\cdot q=0.
\]
For a scalar function $f:V_m\to\mathbb R$, define the row functional
\begin{equation}\label{eq:delta}
 \Delta_{a,p,q}f
 =f(a)+f(a+p+q)-f(a+p)-f(a+q).
\end{equation}
Reordering the two diagonals changes at most the sign of this row and does
not change its singular values.

Let $\mathcal R_m$ be a set of distinct elastic rectangles.  The signed
matrix $C_m\in\{-1,0,1\}^{|\mathcal R_m|\times N_m}$ has one row per
rectangle, with the signs in \eqref{eq:delta}.  Each row has four nonzero
entries.  Its incidence degree at a vertex $v$ is
\[
 d_v=|\{R\in\mathcal R_m:v\in R\}|,
 \qquad D_m=\operatorname{diag}(d_v:v\in V_m).
\]
Throughout the normalized discussion we require $d_v>0$ for every vertex.
If a vertex is isolated, $D_m^{-1/2}$ is not defined in the intended
normalization; moreover, the associated coordinate vector supplies an
extra unconstrained direction in the unnormalized row matrix.

\subsection{Collision modes and the canonical gap}

Let
\begin{equation}\label{eq:K}
 \mathcal K_m=\operatorname{span}\{1,x,y,z,x^2+y^2+z^2\}
 \subset\mathbb R^{V_m}.
\end{equation}
Direct substitution in \eqref{eq:delta} shows that $C_mk=0$ for every
$k\in\mathcal K_m$.  The five functions are independent for $m\geq1$.
For instance, evaluating them at
\[
 (0,0,0),\ (1,0,0),\ (-1,0,0),\ (0,1,0),\ (0,0,1)
\]
gives a $5\times5$ matrix whose determinant has absolute value $2$.

We use the degree-normalized signed operator
\begin{equation}\label{eq:A}
 A_m=\frac12 C_mD_m^{-1/2}.
\end{equation}
The factor $1/2$ fixes the scale of a four-point row.  In particular, the
associated combinatorial convention is $4A_m^*A_m$.  The known nullspace in
normalized coordinates is $D_m^{1/2}\mathcal K_m$.

\begin{definition}[Canonical quotient gap]\label{def:gap}
The normalized gap is
\begin{equation}\label{eq:gap}
 \gamma_m=\min\bigl\{\norm{A_mX}_2:
 \norm{X}_2=1,\ X\perp D_m^{1/2}\mathcal K_m\bigr\}.
\end{equation}
Equivalently,
\begin{equation}\label{eq:rayleigh}
 \gamma_m^2=\frac14\min_{f\notin\mathcal K_m}
 \frac{\sum_{R\in\mathcal R_m}|\Delta_Rf|^2}
 {\min_{h\in\mathcal K_m}\sum_{v\in V_m}d_v|f(v)-h(v)|^2}.
\end{equation}
\end{definition}

The quotient in \eqref{eq:gap} removes only the five known modes.  Thus an
additional kernel vector makes $\gamma_m=0$.  Conversely, the exact identity
$\ker_{\mathbb Q}C_m=\mathcal K_m$ for a finite rational matrix does not
supply a scale-independent lower bound on $\gamma_m$.

\section{The weighted-span obstruction}\label{sec:theorem}

We now give the rigorous result of the paper.  Vector lengths are Euclidean.
For each rectangle $R$, choose one representation \eqref{eq:rectangle} and
write its side vectors as $p_R,q_R$.

\begin{proposition}[Weighted-span test]\label{prop:weighted-span}
Let $m\geq1$ and suppose every vertex has positive incidence degree.  Then
the gap in Definition~\ref{def:gap} satisfies
\begin{equation}\label{eq:weighted-chain}
 \gamma_m^2\leq
 G_m^{-1}\sum_{R\in\mathcal R_m}(p_{R,x}q_{R,x}-p_{R,y}q_{R,y})^2
 \leq G_m^{-1}\sum_{R\in\mathcal R_m}|p_R|^2|q_R|^2,
\end{equation}
where
\[
 G_m=\sum_{(x,y,z)\in V_m}(x^2-y^2)^2.
\]
Writing $S_j=\sum_{t=-m}^m t^j$, one has
\begin{align}
 G_m
 &=2S_0(S_0S_4-S_2^2)\label{eq:G-moment}\\
 &=\frac{2(2m+1)^3m(m+1)(4m^2+4m-3)}{45}.\label{eq:G-closed}
\end{align}
\end{proposition}

\begin{proof}
Set
\[
 g(x,y,z)=x^2-y^2.
\]
Cubic symmetry makes $g$ orthogonal, in the unweighted counting inner
product on $V_m$, to all of $\mathcal K_m$.  Indeed, exchanging $x$ and $y$
handles the constant and radial quadratic modes, while sign symmetry handles
the three linear modes.

Use the degree-weighted inner product
\[
 \ip{u}{w}_d=\sum_{v\in V_m}d_vu(v)w(v),
 \qquad \norm{u}_d^2=\ip{u}{u}_d.
\]
Let $h\in\mathcal K_m$ be the degree-weighted orthogonal projection of $g$
onto $\mathcal K_m$, and put $f=g-h$.  Then
$X=D_m^{1/2}f$ is orthogonal to $D_m^{1/2}\mathcal K_m$.  Also,
$C_mh=0$, so $C_mf=C_mg$.  The denominator has the lower bound
\begin{align}
 \norm{X}_2^2
 &=\norm{f}_d^2
 =\min_{k\in\mathcal K_m}\sum_v d_v|g(v)-k(v)|^2\notag\\
 &\geq\min_{k\in\mathcal K_m}\sum_v|g(v)-k(v)|^2
 =\norm{g}_2^2=G_m.\label{eq:den-bound}
\end{align}
Here the inequality uses $d_v\geq1$, and the final equality uses the
unweighted orthogonality of $g$ to $\mathcal K_m$.  In particular, $X$ is
nonzero.

For a rectangle represented by $(a,p,q)$, direct expansion gives
\begin{equation}\label{eq:delta-g}
 \Delta_{a,p,q}g=2(p_xq_x-p_yq_y).
\end{equation}
Using $A_mX=\tfrac12C_mf=\tfrac12C_mg$ and \eqref{eq:delta-g}, we obtain
\begin{equation}\label{eq:num-bound}
 \norm{A_mX}_2^2
 =\frac14\sum_{R\in\mathcal R_m}|\Delta_Rg|^2
 =\sum_{R\in\mathcal R_m}(p_{R,x}q_{R,x}-p_{R,y}q_{R,y})^2.
\end{equation}
The Rayleigh quotient of this admissible $X$, together with
\eqref{eq:den-bound}, proves the first inequality in
\eqref{eq:weighted-chain}.  The second follows from Cauchy--Schwarz:
\[
 |p_xq_x-p_yq_y|
 \leq\sqrt{p_x^2+p_y^2}\sqrt{q_x^2+q_y^2}
 \leq |p|\,|q|.
\]

It remains to evaluate $G_m$.  Expanding and summing first in $x,y$ and then
in $z$ yields
\[
 G_m=S_0\bigl(2S_0S_4-2S_2^2\bigr),
\]
which is \eqref{eq:G-moment}.  The standard finite sums are
\[
 S_0=2m+1,\qquad
 S_2=\frac{m(m+1)(2m+1)}{3},\qquad
 S_4=\frac{m(m+1)(2m+1)(3m^2+3m-1)}{15}.
\]
Substitution gives \eqref{eq:G-closed}.
\end{proof}

\begin{corollary}[Sublinear spans give no uniform gap]\label{cor:sublinear}
Let $(\mathcal R_m)$ be a family for unbounded $m$.  Suppose every vertex has
positive degree, $d_v\leq D$ with $D$ independent of $m$, and every rectangle
has $|p_R|,|q_R|\leq L_m$.  Then
\begin{equation}\label{eq:sublinear}
 \gamma_m\leq
 \frac{L_m^2}{2}
 \sqrt{\frac{45D}{2m(m+1)(4m^2+4m-3)}}
 =O_D\bigl((L_m/m)^2\bigr).
\end{equation}
In particular, $L_m=o(m)$ implies $\gamma_m\to0$.  Fixed-span families have
$L_m=O(1)$ and recover the older $O_D(m^{-2})$ local obstruction.  More
generally, any sublinear-span family fails to be a uniform expander under
\eqref{eq:A}--\eqref{eq:gap}.  The only span regime not excluded by this
test is diameter scale, $L_m=\Theta(m)$.
\end{corollary}

\begin{proof}
Every row has four distinct vertices, so
$4|\mathcal R_m|=\sum_vd_v\leq D(2m+1)^3$.  Combining this count with
Proposition~\ref{prop:weighted-span} and $|p_R|,|q_R|\leq L_m$ gives
\[
 \gamma_m^2\leq \frac{D(2m+1)^3L_m^4}{4G_m}.
\]
Substituting \eqref{eq:G-closed} gives \eqref{eq:sublinear}.
\end{proof}

\begin{corollary}[Macroscopic-span necessary condition]\label{cor:macro}
Fix constants $D<\infty$ and $c>0$.  Let
\[
 \rho=\min\{1,c/\sqrt{135D}\},
 \qquad
 \eta=\min\{1,c^2/(540D)\}.
\]
If a family on $V_m$ has $1\leq d_v\leq D$ and $\gamma_m\geq c$, then at
least an $\eta$ fraction of its rectangles simultaneously satisfy
\[
 |p_R|\geq\rho m,\qquad |q_R|\geq\rho m.
\]
Thus a scale-independent gap forces a positive density of quantitatively
macroscopic, diameter-scale rectangles.  This condition is necessary only;
it is not sufficient for a gap.
\end{corollary}

\begin{proof}
From \eqref{eq:G-closed}, $G_m\geq(4/45)(2m+1)^3m^4$.  Also
$|\mathcal R_m|\leq D(2m+1)^3/4$.  In the cube $V_m$, every side vector has
length at most $2\sqrt3\,m$, so $|p_R|^2|q_R|^2\leq144m^4$.

Call a rectangle bad if at least one of its two side lengths is $<\rho m$.
For a bad rectangle, $|p_R|^2|q_R|^2\leq12\rho^2m^4$.  Therefore the total
bad contribution is at most
$3D\rho^2(2m+1)^3m^4\leq(c^2/45)(2m+1)^3m^4$.  Since $\gamma_m\geq c$,
Proposition~\ref{prop:weighted-span} forces
\[
 \sum_R|p_R|^2|q_R|^2\geq c^2G_m
 \geq (4c^2/45)(2m+1)^3m^4.
\]
The good rectangles therefore contribute at least
$(c^2/15)(2m+1)^3m^4$.  Each contributes at most $144m^4$, so their number is
at least $c^2(2m+1)^3/2160$.  Dividing by
$|\mathcal R_m|\leq D(2m+1)^3/4$ gives a fraction at least
$c^2/(540D)$, and hence at least $\eta$.
\end{proof}

\begin{remark}[What remains open]\label{rem:nonlocal}
Corollary~\ref{cor:sublinear} rules out fixed, local, block-glued, and all
sublinear-span bounded-degree families under the stated normalization.
Corollary~\ref{cor:macro} says any positive answer to the nonlocal problem
must place a positive fraction of rows in the diameter-scale regime.  The
paper does not prove that such a macroscopic family exists, nor that every
macroscopic family has a gap.
\end{remark}

\section{Canonical finite pools and exact kernel certificates}
\label{sec:certificates}

The finite experiment used all nonperiodic rectangles in $V_m$ as its
candidate pool.  An unordered diagonal is a pair $\{u,w\}\subset V_m$.
Diagonals were grouped by
\[
 (u+w,\ |u-w|^2).
\]
Two distinct diagonals in one group have a common midpoint and equal length;
their four endpoints form an elastic rectangle.  Sorting the pairs and
assigning $+,+$ to the earlier diagonal and $-,-$ to the later one produces a
canonical row.  Pairing every two distinct diagonals in every group gives the
complete finite pool under this convention.  Shared-vertex degeneracy cannot
occur between two distinct pairs in one group: equal sums would force the
other endpoints to coincide.  Periodic boundaries were not used, because
reducing coordinates modulo a period would not preserve the quadratic energy
mode on the displayed representatives.

The designated finite instances were selected by a residual-greedy branch
from these complete pools.  The branch begins with a degree-capped balanced
selection and then adds rows having large residual in the current numerical
quotient computation.  This is a search heuristic, not an infinite recursive
construction and not an exhaustive optimization.

For each designated candidate, the exact certificate performs the following
checks.
\begin{enumerate}
\item Every row is unique, belongs to the canonical pool, and satisfies the
integer momentum, energy, and orthogonality identities.
\item All five integer mode columns in \eqref{eq:K} have zero row residual,
and their anchor determinant has absolute value $2$.  Therefore
$\operatorname{rank}_{\mathbb Q}C_m\leq N_m-5$.
\item The certificate lists row and column indices of an
$(N_m-5)\times(N_m-5)$ integer minor whose determinant is $1$ modulo $2$.
The integer determinant is therefore nonzero, so
$\operatorname{rank}_{\mathbb Q}C_m\geq N_m-5$.
\end{enumerate}
It follows exactly, for that finite matrix, that
\[
 \operatorname{rank}_{\mathbb Q}C_m=N_m-5,
 \qquad \ker_{\mathbb Q}C_m=\mathcal K_m.
\]
Ranks over the two large primes $2147483647$ and $2147483629$ were also
computed and equal $N_m-5$.  Those prime computations are independent
screens; the odd maximal minor, not the prime screens, is the rational-rank
proof.

\section{Four finite full-pool instances}\label{sec:finite}

Table~\ref{tab:candidates} records the four designated residual-greedy
instances.  The vertex and row counts, degree range, rational rank, and
kernel dimension are exact finite data (E).  The final column is a
floating-point normalized gap (N), freshly recomputed from the raw JSON row
list with the normalization \eqref{eq:A}.  Digits shown are rounded and are
not certified intervals.

\begin{table}[ht]
\centering
\small
\caption{Finite residual-greedy instances.  Columns through
$\dim\ker_{\mathbb Q}$ are exact finite data; $\gamma_m$ is numerical.}
\label{tab:candidates}
\begin{tabular}{@{}rrrrrrr@{}}
\toprule
$m$ & $N_m$ & $|\mathcal R_m|$ & degree range &
$\operatorname{rank}_{\mathbb Q}C_m$ & $\dim\ker_{\mathbb Q}C_m$ &
normalized $\gamma_m$ \\
\midrule
2 & 125  & 188  & $4$--$10$ & 120  & 5 & 0.137567327104 \\
3 & 343  & 515  & $3$--$10$ & 338  & 5 & 0.108836005898 \\
4 & 729  & 1094 & $3$--$10$ & 724  & 5 & 0.107304440463 \\
5 & 1331 & 1997 & $3$--$10$ & 1326 & 5 & 0.104461520420 \\
\bottomrule
\end{tabular}
\end{table}

The independent numerical recomputation rebuilt $C_m$, $D_m$, $A_m$, and the
weighted constraint basis directly from the four row arrays.  A fresh
constrained eigensolve gave, respectively,
\[
 0.13756732710431338,\quad
 0.10883600589808418,\quad
 0.10730444046255100,\quad
 0.10446152042018168.
\]
The fresh quotient-operator residuals were all below $6\times10^{-10}$, and
no solver warning was emitted in those runs.  These checks support the
reported finite approximations.  They do not turn the gaps into exact
algebraic values.

Two non-implications are essential.  First, exact five-dimensional kernels
at four sizes do not imply exact kernels at later sizes.  Second, the observed
numbers do not imply that $\liminf_{m\to\infty}\gamma_m>0$.  The candidates
were chosen from full, nonperiodic pools without a fixed side-length cutoff,
so Corollary~\ref{cor:sublinear} does not itself rule them out; conversely, the
four candidates do not define or certify an infinite nonlocal family.

\section{Validation and reproducibility}\label{sec:validation}

The evidence package separates symbolic, exact finite, and numerical checks.
The following layers were used.
\begin{enumerate}
\item \textbf{Formula and pool checks.}  Integer arithmetic verifies every
canonical row, hashes the ordered row arrays with domain-separated SHA-256,
and rejects duplicate rows.
\item \textbf{Exact kernel checks.}  The odd-minor procedure described in
Section~\ref{sec:certificates} was rerun independently for all four designated
candidates.  Fresh proof payloads agree with the packaged payloads after
excluding run metadata.
\item \textbf{Raw-gap recomputation.}  The normalized operators and
constraint bases were reconstructed from the raw row arrays rather than from
summary tables.  Stored numerical vectors and fresh constrained eigenvectors
were checked for norm, constraint residual, operator residual, and relative
gap agreement.
\item \textbf{Theorem arithmetic checks.}  Exact integer summation checks
\eqref{eq:G-moment} and the five orthogonality identities for $m=1,\ldots,10$.
On four finite local block candidates, a separate script checks the identity
\eqref{eq:delta-g}, the side-length bound, the projected Rayleigh quotient,
and the inequality chain in Proposition~\ref{prop:weighted-span}.
\item \textbf{Manuscript-local checks.}  The draft package contains a source
and claim checker, build logs, and a source matrix mapping each strong claim
to an exact local artifact or to one of the two verified citations.
\end{enumerate}
The finite checks in items 3 and 4 are validation aids.  The general theorem
depends on the proof in Section~\ref{sec:theorem}, not on extrapolation from
those checks.

\section{Literature-boundary checks}\label{sec:literature}

This revision made two targeted public literature checks.  First, in discrete
velocity models for the Boltzmann equation, the standard physical invariants
are mass, momentum, and energy; extra invariants are called spurious or
nonphysical, and DVMs without them are called normal
\cite{BobylevCercignani1999DVM,BernhoffVinerean2016Mixtures}.  This is close
in vocabulary to the finite statement $\ker C=\operatorname{span}\{1,v_x,v_y,
v_z,|v|^2\}$, but the checked sources concern kinetic DVM construction and
mixtures, not this paper's finite cubic-lattice rectangle-frame quotient gap.

Second, the locality obstruction is consistent with standard finite-state
Poincar\'e and spectral-gap intuition: bounded-range walks on long one-
dimensional or lattice-like sets have gaps that decay with the diameter; for
example, the simple random walk on the $n$-cycle has spectral gap
$1-\cos(2\pi/n)$ \cite{GoelMontenegroTetali2006SpectralProfile}.  This paper's
proof is nevertheless a separate test-function argument for the specific
four-point collision operator, five-mode quotient, $1/2$ normalization, and
constants above.

\begin{center}
\small
\begin{tabular}{@{}p{0.20\textwidth}p{0.24\textwidth}p{0.18\textwidth}p{0.22\textwidth}@{}}
\toprule
Claim & closest prior result checked & same/different & safe wording \\
\midrule
Five physical collision modes & normal DVMs have no spurious collision
invariants beyond mass, momentum and energy & same invariant vocabulary;
different finite matrix and no gap theorem & exact finite kernels here are
normal-DVM-like evidence, not a novelty or priority claim \\
Locality obstructs uniform gap & cycle/lattice spectral-gap decay under
local moves & same Poincar\'e intuition; different operator and quotient &
the paper proves a source-bound weighted-span obstruction for this model \\
Finite $m=2,\ldots,5$ certificates & no direct finite cubic elastic-rectangle
coverage found in the targeted check & related only by invariant language &
certificates retain independent finite meaning; no uniform-gap evidence \\
\bottomrule
\end{tabular}
\end{center}

\section{The remaining nonlocal problem}\label{sec:open}

A natural grid form of the unresolved question is the following.

\begin{quote}
\textbf{Open nonlocal rectangle-frame proposition.}
Do there exist absolute constants $D<\infty$ and $c>0$, and elastic-rectangle
families on velocity grids $V_m$ for unbounded $m$, such that every vertex has
degree between $1$ and $D$,
$\ker_{\mathbb Q}C_m=\mathcal K_m$, and the normalized gap satisfies
$\gamma_m\geq c$, when rectangle side lengths are allowed to grow with $m$?
\end{quote}

Corollary~\ref{cor:macro} identifies a necessary escape from locality: any
positive answer must use a positive fraction of rectangles whose two side
lengths are both bounded below by a fixed multiple of $m$ (or must change
another stated hypothesis).  This is only a necessary condition.
The four values in Table~\ref{tab:candidates} neither prove nor disprove the
open proposition.  The broader existence question on other growing rational
or integer velocity sets is likewise not settled by this draft.

\section{Limitations and next technical step}\label{sec:limits}

The main limitations are structural rather than cosmetic.
\begin{itemize}
\item The theorem is one-sided and conditional on positive bounded degree.  It
rules out sublinear span but supplies no lower bound and no sharpness result.
\item The exact certificates concern four finite matrices only.  They do not
provide a recursive construction.
\item The spectral values are floating-point estimates, not interval or
symbolic eigenvalue certificates.
\item The computational search and attacks are finite.  They are not claimed
to cover all selection mechanisms, all nonlocal patterns, or all velocity
sets.
\item The two targeted literature checks in Section~\ref{sec:literature} are
not a systematic priority review of rectangle-frame Poincar\'e inequalities.
This paper therefore makes no novelty, priority, ``current best,'' or
no-precedent claim.
\end{itemize}

The exact next mathematical step is to prove or refute a scale-free
Poincar\'e inequality for an explicit nonlocal bounded-degree rectangle
family, with the five collision modes as the full rational kernel.  A
positive route needs a recursive construction and a uniform lower bound; a
negative route needs a test function or decomposition that remains effective
in the quantitatively macroscopic side-length regime.  Extending the finite
table alone is not a substitute for either argument.  Before formal submission, a systematic
literature review and independent mathematical review are also required.

\section{Conclusion}

The rigorous conclusion is narrow and explicit: on $V_m$, the weighted-span
test rules out uniformly bounded-degree families whose rectangle side lengths
are $o(m)$.  A scale-independent gap would require a positive fraction of
diameter-scale rectangles.  Exact finite certificates at $m=2,3,4,5$ establish
five-dimensional rational kernels for four selected matrices, and numerical
recomputation supports their displayed finite gaps.  Those observations do
not settle the unrestricted nonlocal proposition, which remains open.

\section*{AI authorship, assistance, and responsibility disclosure}

Every sentence of manuscript prose and the complete manuscript structure were
generated by LingTai AI under Runyuan Wang's direction and authorization.
LingTai AI also performed source organization, mathematical and computational
analysis assistance, and code-assisted validation.  The original conjectural
agenda and experimental direction came from Runyuan Wang.  AI production is
not evidence for any claim: the stated proofs, source-bound checks, finite
certificates, and disclosed numerical computations must support the paper on
their own.  Independent human review remains mandatory before external use or
submission.  Runyuan Wang retains final scientific judgment and responsibility
for any public or submitted version.  This work is shared purely for personal
enjoyment and playful exploration with AI.  If anything is incorrect, please
excuse it; criticism and corrections are warmly welcomed.

\bibliographystyle{plain}
\bibliography{references}

\end{document}
